\section{A published large-support estimate sharply bounds the residual symplectic rank}
\label{sec:symplectic-three-sharp}

\paragraph{Status and relation to the preceding proof.}
This separate addendum leaves D-HOS-13--15 unchanged. Its only new
published input is the explicit estimate and support refinement of
Larsen--Shalev--Tiep. All constants and the passage through the perfect
cover are recorded below. The entire argument and invoked published proof have passed independent
whole-proof reviews by /root and /root/imprimitive; the separate exact
arithmetic also passes. It does not assert existence at any surviving rank.
As before, $\nu(T)$ is the codimension of the largest complex eigenspace
of $T$, and $n$ denotes half the natural dimension of
$S=\operatorname{Sp}_{2n}(3)$.

\begin{theorem}[D-HOS-16: a finite residual rank interval]
\label{D-HOS-16}
Put
\[
 c_*=(19/20)^9=\frac{322687697779}{512000000000},\qquad
 \delta_*:=\frac{1-c_*}{2}
          =\frac{189312302221}{1024000000000}.
 \tag{HOS3S.1}
\]
For every $n\ge284$, every nontrivial irreducible ordinary complex
$S$-module $V$ of degree $a$, and every noncentral $x\in S$,
\[
                           \nu(x|V)\ge\delta_*a.
 \tag{HOS3S.2}
\]
Consequently a finite inducing image with an irreducibly acting normal
central product of quasisimple factors whose simple quotients are
$\operatorname{PSp}_{2n_i}(3)$, $n_i\ge284$, and whose factor modules
are all nonselfdual cannot have period-image dimension nine for its
actual six-dimensional multiplicity variation. Arbitrary central
quotients, scalar extensions, different such ranks, and compatible
factor permutations are allowed.

In the transitive characteristic-three symplectic-core case of the
main imaginary-quadratic sixfold problem, D-HOS-12 and this theorem
therefore reduce the remaining case with $n\ge4$ to
\[
                                  4\le n\le283.
 \tag{HOS3S.3}
\]
\end{theorem}

\begin{proof}
We first specify the primary estimate, including the distinction
between its two support conventions. Larsen--Shalev--Tiep,
\cite[Definition~4.1.1 and (4.1.1), p.~1907; Theorem~4.3.6,
pp.~1917--1918, (4.3.2)]{HOS:LST2011}, treat a simple simply
connected classical algebraic group $\mathbf G$, its finite group
$\mathbf G(\mathbf F_q)$, and every nontrivial ordinary irreducible
character $\chi$ of that finite group. If the support of $h$ is at
least $N\ge567$, their explicit estimate is
\[
 \frac{|\chi(h)|}{\chi(1)}
       <\gamma_q^{\sqrt{N/7}},\qquad \gamma_3=19/20.
 \tag{HOS3S.4}
\]
Here the usual support is the codimension of the largest eigenspace
over the algebraic closure in the natural representation. The
radical covers $N/7$: the complete proof, part~4 on p.~1922,
explicitly chooses $\alpha=\sqrt{1/7}$ in its exponent
$\alpha\sqrt N$. Thus $N=567=7\cdot9^2$ gives the constant
$c_*$ in (HOS3S.1).

The same authors then state, in Remark~4.3.7, p.~1923, that their
conclusion also holds with support restricted to the allowed scalar
set. For the symplectic group this set is $\mu_2=\{1,-1\}$.
We may therefore use (HOS3S.4) with
\[
 s_{\pm}(h):=\min_{\lambda\in\{1,-1\}}
             \operatorname{codim}\ker(h-\lambda I).
 \tag{HOS3S.5}
\]
This is an explicit published refinement, not an identification
of $s_{\pm}$ with the usual algebraic-closure support.

Apply these statements to the algebraic group
$\mathbf G=\operatorname{Sp}_{2n}$, whose rational points are
the full perfect cover $S$, including its central involution.
The theorem does not require $\chi$ to be faithful. Thus it applies
unchanged to a character inflated from any represented central
quotient; no character-ratio assertion about an unspecified quotient
is being substituted for the published statement.

Take a power of $x$ of prime order $\ell$ modulo $Z(S)$, and
change its central sign as necessary, calling the result $h$.
The projective kernel of a nontrivial irreducible $S$-module is
exactly $Z(S)$, as in D-HOS-14. Taking powers and multiplying by
a scalar cannot increase projective support. If $\ell$ is odd,
choose the sign so that $h$ has actual order $\ell$; if $\ell=2$,
then $h^2=I$ or $-I$.

All but one family of cases have the uniform estimates already
proved in D-HOS-13 and the elementary part of D-HOS-14. For
actual order three the bound is $2a/9$. If $h^2=I$, choose its
central sign so that the smaller nonzero symplectic eigenspace is
fixed, and remove a fixed hyperbolic plane. The quotient moving
rank is at least $n\ge4$, so the phase-tolerant Heisenberg estimate
(HOS3C.8) gives at least $(1/4)(1-3^{-2})a=2a/9$.
For odd $\ell\ge5$ with a natural eigenvalue in $\{1,-1\}$,
the actual odd-order representative has a fixed symplectic plane;
the quotient moving rank is at least two. The same estimate gives
at least $(1/2)(4/5)(2/3)a=4a/15$.

The residual cases are semisimple and have neither $1$ nor $-1$
as a natural eigenvalue, including $h^2=-I$. The usual
algebraic-closure support is at least $n$, because each minimal
polynomial factor over $\mathbf F_3$ has degree at least two.
Thus the unrefined theorem alone already gives the argument for
$n\ge567$. Using the explicit published refinement (HOS3S.5),
however, we have exactly $s_{\pm}(h)=2n$. For $n\ge284$ this
is at least $568>567$, so (HOS3S.4) with $N=567$ gives
$|\chi(h)|/a<c_*$. Rotating an eigenvalue of maximum multiplicity
to $1$ and applying the triangle inequality gives
$|\chi(h)|\ge a-2\nu(h|V)$. Hence $\nu(h|V)>\delta_*a$.
The exact inequalities
\[
                 \frac16<\delta_*<\frac29<\frac14
 \tag{HOS3S.6}
\]
show that the preceding elementary cases also imply (HOS3S.2).
This proves the assertion for every noncentral $x$.

For the source consequence we repeat the normalizer passage so
its scope is explicit. In the stated rank range the universal
perfect cover of each simple factor is $\operatorname{Sp}_{2n_i}(3)$,
and its only nontrivial outer class is the nonsquare-multiplier
diagonal class, by the cover and automorphism inputs already
checked in D-HOS-15. Vinroot's anti-symplectic factorization
implies that this class sends each character to its complex
conjugate. It therefore does not stabilize a nonselfdual factor
module. Every compatible factor normalizer is consequently a
scalar times an inner operator. This remains true after passing
to a represented central quotient, since its automorphisms lift
to the perfect cover and Schur's lemma identifies the remaining
intertwiner scalar.

On the full tensor space of dimension $e$, a factor-preserving
nonscalar operator has a nonscalar factor, so (HOS3S.2) and
convolution of eigenvalue multiplicities give support at least
$\delta_*e$. An operator permuting factors nontrivially has
support at least $e/4$ by the weighted tensor-cycle lemma, which
is stronger. A nonidentity scalar moves the entire space.
Thus every genuine inertia moves at least $\delta_*e$ dimensions.

The actual intermediate-cover formula is
$6=2e(h-1)+D$. The elementary degree bound D-HOS-02 gives
$e\ge3^{283}$, so genus at least two is impossible. In genus
one, $D=6$ permits at most six genuine marks and hence period
rank at most six. Rank nine would force genus zero, at least
twelve genuine inertia points, and $D=2e+6$. It would follow that
\[
 e\le\frac{6}{12\delta_*-2}
   =\frac{1536000000000}{55936906663}<28,
 \tag{HOS3S.7}
\]
contradicting $e\ge3^{283}$. This is a bound on the actual
multiplicity variation and does not identify the inducing
character field with the original imaginary quadratic field.

Finally selfduality is invariant under automorphic pullback and
module isomorphism. In a transitive core all factors therefore
have the same selfduality status. The selfdual case has already
been excluded by D-HOS-12 in the main imaginary-character-field
setting. The nonselfdual case with $n\ge284$ is excluded above.
This proves (HOS3S.3); it makes no assertion about the existence
or exclusion of the remaining nonselfdual ranks $4$ through $283$.
\end{proof}

\paragraph{Certificate boundary.}
The accompanying rational certificate verifies the displayed
constants, support thresholds, and the source degree contradiction.
It does not enumerate groups or certify the published character
theorem. The latter is used at the exact stated primary locators,
and the full proof of Theorem~4.3.6 and its subsequent remark were
read in the published version. The new addendum is independent
of the very small explicit exponent used in D-HOS-14.
